% ==============================================================================
% Paper 86 (EN): The Rank of Elliptic Curves: Boundedness and Distribution
% Author: Michael Fuhrmann
% Project: Specialist Mathematics Project, Build 167
% Date: 2026-03-12
% ==============================================================================

\documentclass[12pt,a4paper]{amsart}

\usepackage{amsmath,amssymb,amsthm,mathtools,enumitem,booktabs,array,hyperref}
\usepackage[utf8]{inputenc}
\usepackage[T1]{fontenc}
\usepackage{geometry}
\geometry{margin=2.5cm}

% --------------------------------------------------------------------------
% Theorem environments
% --------------------------------------------------------------------------
\newtheorem{theorem}{Theorem}[section]
\newtheorem{lemma}[theorem]{Lemma}
\newtheorem{corollary}[theorem]{Corollary}
\newtheorem{proposition}[theorem]{Proposition}
\theoremstyle{definition}
\newtheorem{definition}[theorem]{Definition}
\newtheorem{example}[theorem]{Example}
\theoremstyle{remark}
\newtheorem{remark}[theorem]{Remark}
\newtheorem{conjecture}[theorem]{Conjecture}
\newtheorem{algorithm}[theorem]{Algorithm}

% --------------------------------------------------------------------------
% Macros
% --------------------------------------------------------------------------
\newcommand{\Z}{\mathbb{Z}}
\newcommand{\Q}{\mathbb{Q}}
\newcommand{\R}{\mathbb{R}}
\newcommand{\C}{\mathbb{C}}
\newcommand{\F}{\mathbb{F}}
\newcommand{\N}{\mathbb{N}}
\newcommand{\Qp}{\mathbb{Q}_p}
\newcommand{\Gal}{\mathrm{Gal}}
\newcommand{\Sel}{\mathrm{Sel}}
\DeclareMathOperator{\Sha}{\mathfrak{Sh}}
\newcommand{\rank}{\mathrm{rank}}
\newcommand{\ord}{\mathrm{ord}}
\newcommand{\Frob}{\mathrm{Frob}}
\newcommand{\Qbar}{\overline{\Q}}
\newcommand{\Prob}{\mathrm{Prob}}
\DeclareMathOperator{\Nm}{Nm}
\newcommand{\LE}{L(E,s)}

\hypersetup{
  colorlinks=true,
  linkcolor=blue,
  citecolor=blue,
  urlcolor=blue,
  pdftitle={The Rank of Elliptic Curves: Boundedness and Distribution},
  pdfauthor={Michael Fuhrmann}
}

% ==============================================================================
\begin{document}
% ==============================================================================

\title{The Rank of Elliptic Curves:\\
  Boundedness and Distribution}

\author{Michael Fuhrmann}
\address{Specialist Mathematics Project, Build 167}
\email{specialist-maths@project.local}
\date{2026-03-12}

\begin{abstract}
The Mordell-Weil theorem asserts that for an elliptic curve $E/\Q$, the group $E(\Q)$
of rational points is a finitely generated abelian group:
$E(\Q) \cong \Z^r \oplus E(\Q)_{\mathrm{tors}}$.
The \emph{rank} $r = r(E)$ measures the ``size'' of the rational points and is one of
the most enigmatic invariants in number theory.
We survey the theory of the rank: the Mordell-Weil theorem (via descent), the Selmer
group and Tate-Shafarevich group $\Sha(E/\Q)$, the analytic rank via the
$L$-function $L(E,s)$ and the Birch and Swinnerton-Dyer conjecture (BSD),
records (Elkies' curve of rank $\geq 29$), the Goldfeld conjecture on average rank $1/2$,
the boundedness conjecture (rank bounded by an absolute constant),
Bhargava-Shankar's theorem (average rank $\leq 0.885$, positive proportion of rank $0$
and rank $1$ curves), the Poonen-Rains probabilistic model, and the parity conjecture.
The boundedness conjecture and Goldfeld's conjecture both remain open.
\end{abstract}

\maketitle

\tableofcontents

% ==============================================================================
\section{Introduction}
% ==============================================================================

An \emph{elliptic curve} over $\Q$ is a smooth projective curve of genus $1$ with a
distinguished rational point $O$ (the ``point at infinity''), given in Weierstrass form
\[
  E\colon y^2 + a_1 xy + a_3 y = x^3 + a_2 x^2 + a_4 x + a_6, \quad a_i \in \Z,
\]
with non-zero discriminant. The set $E(\Q)$ of rational points forms an abelian group
under the chord-and-tangent law.

\begin{theorem}[Mordell 1922, Weil 1928 \cite{Mordell1922,Weil1928}]
  For any elliptic curve $E/\Q$, the group $E(\Q)$ is finitely generated:
  \[
    E(\Q) \;\cong\; \Z^r \oplus E(\Q)_{\mathrm{tors}},
  \]
  where $r = r(E) \geq 0$ is the \emph{rank} of $E$ and $E(\Q)_{\mathrm{tors}}$ is a
  finite abelian group.
\end{theorem}

The rank $r(E)$ is the central object of study in this paper. Key questions:
\begin{itemize}
  \item How large can $r(E)$ be? (Unboundedness / boundedness)
  \item What is the ``typical'' rank? (Distribution / average rank)
  \item Is there an algorithm to compute $r(E)$? (Effective BSD)
\end{itemize}

% ==============================================================================
\section{The Mordell-Weil Theorem: Proof Sketch via Descent}\label{sec:mordell_weil}
% ==============================================================================

The proof of the Mordell-Weil theorem proceeds in two steps.

\subsection{Weak Mordell-Weil theorem}

\begin{theorem}[Weak Mordell-Weil]
  For every $n \geq 2$, the quotient group $E(\Q)/nE(\Q)$ is finite.
\end{theorem}

\begin{proof}[Sketch]
  The key is a height argument. Fix $n=2$. The $2$-isogeny sequence gives a map
  $E(\Q) \to H^1(\Gal(\Qbar/\Q), E[2])$, landing in a finite Selmer group.
  The finiteness of the image comes from the finiteness of $S$-units and Kummer theory.
\end{proof}

\subsection{Descent and heights}

\begin{definition}[Weil height machine]
  The \emph{canonical (Néron-Tate) height} $\hat{h}\colon E(\Q) \to \R_{\geq 0}$ is a
  positive semi-definite quadratic form on $E(\Q) \otimes \R$ with $\hat{h}(P) = 0$
  iff $P$ is torsion.
\end{definition}

The Mordell-Weil theorem follows: finitely generated is proved by combining weak
Mordell-Weil with the fact that $\{P \in E(\Q) : \hat{h}(P) \leq C\}$ is finite for any $C$.

% ==============================================================================
\section{Torsion Subgroup}\label{sec:torsion}
% ==============================================================================

\begin{theorem}[Mazur, 1977 \cite{Mazur1977}]
  For any elliptic curve $E/\Q$, the torsion subgroup $E(\Q)_{\mathrm{tors}}$ is
  isomorphic to one of the $15$ groups:
  \begin{align*}
    &\Z/n\Z \quad (n = 1,2,\ldots,10,12), \\
    &\Z/2\Z \times \Z/2n\Z \quad (n = 1,2,3,4).
  \end{align*}
\end{theorem}

\begin{remark}
  Every one of these $15$ torsion structures is realized by some elliptic curve over $\Q$.
  The most common are $\Z/1\Z$ (trivial) and $\Z/2\Z$.
\end{remark}

% ==============================================================================
\section{The Selmer Group and Tate-Shafarevich Group}\label{sec:selmer}
% ==============================================================================

\begin{definition}[Selmer group and $\Sha$]
  For a prime $p$ and an elliptic curve $E/\Q$, the short exact sequence
  $0 \to E[p] \to E \xrightarrow{[p]} E \to 0$
  gives (via Galois cohomology) a long exact sequence:
  \[
    0 \;\to\; \frac{E(\Q)}{pE(\Q)} \;\to\; \Sel^{(p)}(E/\Q) \;\to\; \Sha(E/\Q)[p] \;\to\; 0.
  \]
  The \emph{$p$-Selmer group} $\Sel^{(p)}(E/\Q)$ consists of cohomology classes in
  $H^1(\Q, E[p])$ that are locally trivial at all primes $v$:
  \[
    \Sel^{(p)}(E/\Q) = \ker\left(H^1(\Q, E[p]) \to \prod_v H^1(\Q_v, E)\right).
  \]
  The \emph{Tate-Shafarevich group} is
  \[
    \Sha(E/\Q) = \ker\left(H^1(\Q, E) \to \prod_v H^1(\Q_v, E)\right).
  \]
\end{definition}

\begin{conjecture}[Finiteness of $\Sha$]\label{conj:sha}
  For every elliptic curve $E/\Q$, the group $\Sha(E/\Q)$ is finite.
\end{conjecture}

\begin{remark}
  Finiteness of $\Sha$ is proven only in special cases:
  \begin{itemize}
    \item when $\ord_{s=1}L(E,s) = 0$: Kolyvagin \cite{Kolyvagin1988} proved $\rank E(\Q) = 0$ and $\Sha(E/\Q)$ finite (using Euler systems);
    \item when $\ord_{s=1}L(E,s) = 1$: Gross-Zagier \cite{GZ1986} + Kolyvagin proved $\rank E(\Q) = 1$ and $\Sha(E/\Q)$ finite;
    \item for specific families using Heegner points.
  \end{itemize}
  In general, Conjecture~\ref{conj:sha} is open.
\end{remark}

% ==============================================================================
\section{The $L$-Function and the BSD Conjecture}\label{sec:bsd}
% ==============================================================================

\begin{definition}[$L$-function of an elliptic curve]
  For a prime $p$ of good reduction, define the local $L$-factor:
  \[
    L_p(E,s) = \frac{1}{1 - a_p p^{-s} + p^{1-2s}}, \quad
    a_p = p + 1 - \#E(\F_p).
  \]
  The global $L$-function (completed with Euler factors at bad primes) is
  \[
    L(E,s) = \prod_p L_p(E,s),
  \]
  convergent for $\mathrm{Re}(s) > 3/2$ (by the Hasse bound $|a_p| \leq 2\sqrt{p}$),
  and analytically continued to all of $\C$ (by modularity, Wiles et al.\ \cite{BCDT2001}).
\end{definition}

\begin{conjecture}[Birch and Swinnerton-Dyer \cite{BSD1965}]\label{conj:bsd}
  Let $E/\Q$ be an elliptic curve. Then:
  \begin{enumerate}[label=(\roman*)]
    \item $\ord_{s=1} L(E,s) = \rank E(\Q)$ \quad (analytic rank $=$ algebraic rank);
    \item The leading coefficient of $L(E,s)$ at $s=1$ is given by the BSD formula:
    \[
      \lim_{s\to 1} \frac{L(E,s)}{(s-1)^r} =
      \frac{\Omega_E \cdot \mathrm{Reg}(E/\Q) \cdot |\Sha(E/\Q)| \cdot \prod_p c_p}
           {|E(\Q)_{\mathrm{tors}}|^2},
    \]
    where $\Omega_E$ is the real period, $\mathrm{Reg}(E/\Q)$ is the Néron-Tate
    regulator, $c_p$ are the Tamagawa numbers, and $r = \rank E(\Q)$.
  \end{enumerate}
\end{conjecture}

\begin{remark}
  The BSD conjecture is one of the Millennium Prize Problems (Clay Mathematics Institute).
  The best known result is:
\end{remark}

\begin{theorem}[Coates-Wiles \cite{CW1977}, Kolyvagin \cite{Kolyvagin1988},
  Gross-Zagier \cite{GZ1986}]
  Let $E/\Q$ be an elliptic curve with $L(E,1) \neq 0$. Then $\rank E(\Q) = 0$.
  If $L(E,1) = 0$ and $L'(E,1) \neq 0$ (analytic rank $1$), and $E$ is a CM curve or
  $E$ has a Heegner point of infinite order, then $\rank E(\Q) = 1$ and $\Sha(E/\Q)$
  is finite.
\end{theorem}

% ==============================================================================
\section{Records: High Rank Elliptic Curves}\label{sec:records}
% ==============================================================================

\begin{theorem}[Elkies, 2006 \cite{Elkies2006}]
  There exists an elliptic curve over $\Q$ with rank $\geq 29$.
  The current record (as of 2026) is $\rank \geq 29$, held by the curve constructed by
  Noam Elkies.
\end{theorem}

\begin{remark}
  The curve has equation
  \begin{align*}
    y^2 + xy + y = x^3 &- 126846046810060579589x \\
                        &+ 1669662011297869897320473097301,
  \end{align*}
  (approximately; the exact form is long). The $29$ independent generators are known
  explicitly.
\end{remark}

\begin{remark}
  Prior records: Mestre \cite{Mestre1982} (rank $\geq 12$, 1982),
  Fermigier (rank $\geq 19$, 1997),
  Martin-McMillen (rank $\geq 24$, 2000),
  Elkies (rank $\geq 28$, 2006; rank $\geq 29$, 2006).
\end{remark}

% ==============================================================================
\section{The Goldfeld and Boundedness Conjectures}\label{sec:goldfeld}
% ==============================================================================

\begin{conjecture}[Goldfeld \cite{Goldfeld1979}]\label{conj:goldfeld}
  For elliptic curves $E/\Q$ ordered by conductor $N$,
  \[
    \lim_{X\to\infty} \frac{\sum_{N(E) \leq X} \rank E(\Q)}
                           {\#\{E : N(E) \leq X\}}
    \;=\; \frac{1}{2}.
  \]
  Equivalently, $50\%$ of elliptic curves (ordered by conductor or by naive height)
  have rank $0$, and $50\%$ have rank $1$; the density of rank $\geq 2$ curves is $0$.
\end{conjecture}

\begin{remark}
  The $1/2$ average arises from the parity conjecture: roughly half of elliptic curves
  have even analytic rank and half have odd analytic rank (by the root number), and
  $\ord_{s=1} L = 0$ (resp.\ $1$) is the minimum in each parity class.
\end{remark}

\begin{conjecture}[Boundedness of rank]\label{conj:bounded}
  There exists an absolute constant $C > 0$ such that for all elliptic curves $E/\Q$,
  $\rank E(\Q) \leq C$.
\end{conjecture}

\begin{remark}
  The boundedness conjecture is highly controversial: there exist heuristics supporting
  it (Poonen-Rains \cite{PR2012}, Park-Poonen-Voight-Wood \cite{PPVW2019}) as well as
  constructions of curves with arbitrarily large rank (conditionally on BSD, by
  theoretical methods), suggesting the opposite.
\end{remark}

% ==============================================================================
\section{Bhargava-Shankar: Average Rank Bounds}\label{sec:bhargava}
% ==============================================================================

\begin{theorem}[Bhargava-Shankar \cite{BS2015}]\label{thm:bs}
  When elliptic curves $E/\Q$ are ordered by height
  $H(E) = \max(|a|^3, |b|^2)$ for the model $y^2 = x^3 + ax + b$,
  \begin{enumerate}[label=(\roman*)]
    \item The average size of the $2$-Selmer group $\Sel^{(2)}(E/\Q)$ is $3$.
    \item The average rank satisfies $\overline{r} \leq 0.885$.
    \item A positive proportion of elliptic curves have rank $0$.
    \item A positive proportion of elliptic curves have rank $1$.
  \end{enumerate}
\end{theorem}

\begin{remark}
  The bound $\overline{r} \leq 0.885$ comes from:
  $2\overline{r} \leq \overline{\dim_{\F_2} \Sel^{(2)}} = \log_2(3)$
  via $E(\Q)/2E(\Q) \hookrightarrow \Sel^{(2)}(E/\Q)$.
  More precisely, $2 \cdot \overline{r} \leq 3 - 1 = 2$ (subtracting the torsion
  contribution), giving $\overline{r} \leq 1.5$, but refined arguments give $0.885$.
\end{remark}

\begin{theorem}[Bhargava-Shankar \cite{BS2015a}]
  When ordered by height, the average size of the $3$-Selmer group of $E/\Q$ is $4$.
  Combined with the $5$-descent \cite{BS2013}, this implies $\overline{r} \leq 0.885$.
\end{theorem}

\begin{theorem}[Bhargava-Skinner-Zhang \cite{BSZ2014}]
  A positive proportion of elliptic curves $E/\Q$ have $\rank E(\Q) = 0$ and
  $|\Sha(E/\Q)|$ divides $4$.
  A positive proportion have $\rank E(\Q) = 1$ with $|\Sha|$ finite (conditionally on BSD).
\end{theorem}

% ==============================================================================
\section{The Parity Conjecture}\label{sec:parity}
% ==============================================================================

\begin{definition}[Root number]
  The \emph{root number} $w(E) \in \{+1, -1\}$ appears in the functional equation of
  $L(E,s)$:
  \[
    \Lambda(E,s) = w(E) \cdot \Lambda(E, 2-s),
  \]
  where $\Lambda(E,s) = N^{s/2}(2\pi)^{-s}\Gamma(s) L(E,s)$ is the completed
  $L$-function. Here $N$ is the conductor of $E$.
\end{definition}

\begin{conjecture}[Parity conjecture]\label{conj:parity}
  For every elliptic curve $E/\Q$,
  \[
    (-1)^{\rank E(\Q)} \;=\; w(E).
  \]
  Equivalently, $\rank E(\Q) \equiv \ord_{s=1} L(E,s) \pmod{2}$.
\end{conjecture}

\begin{theorem}[Nekovář \cite{Nekovar2009}]
  The parity conjecture holds for elliptic curves over totally real fields
  $F$ for which $E$ is potentially semistable at all primes above $2$
  (under some local conditions at $2$).
\end{theorem}

\begin{remark}
  The root number $w(E)$ depends only on the local behavior of $E$ at bad primes and at
  $\infty$: $w(E) = \prod_v w_v(E)$. Computing $w(E)$ is effectively doable
  (Rohrlich \cite{Rohrlich1993}).
\end{remark}

% ==============================================================================
\section{The Poonen-Rains Model}\label{sec:poonen_rains}
% ==============================================================================

\begin{conjecture}[Poonen-Rains model \cite{PR2012}]\label{conj:pr}
  Let $p$ be prime. The ``probability'' that a ``random'' elliptic curve $E/\Q$ has
  $p$-Selmer rank $\geq r$ is governed by the random matrix model
  \[
    \Prob[\dim_{\F_p} \Sel^{(p)}(E/\Q) \geq r] \;\sim\;
    \prod_{k=1}^\infty (1 + p^{-k})^{-1} \cdot p^{-r(r-1)/2},
  \]
  modulated by the root number.
\end{conjecture}

\begin{remark}
  Under the Poonen-Rains heuristics:
  \begin{itemize}
    \item The average rank is $1/2$ (consistent with Goldfeld's conjecture);
    \item $50\%$ of curves have rank $0$, $50\%$ have rank $1$;
    \item The proportion of rank $\geq 2$ curves is $0$ (density $0$, but infinitely many);
    \item The proportion of rank $\geq r$ decays super-exponentially: $\sim p^{-r^2}$.
  \end{itemize}
\end{remark}

\begin{conjecture}[Park-Poonen-Voight-Wood \cite{PPVW2019}]
  For each fixed $r \geq 0$, a positive proportion of elliptic curves have rank $r$;
  but the proportion of rank $r$ curves tends to $0$ as $r \to \infty$, consistent
  with a bounded rank conjecture.
\end{conjecture}

% ==============================================================================
\section{The Brumer-McGuinness Data}\label{sec:data}
% ==============================================================================

\begin{remark}
  Brumer and McGuinness \cite{BM1990} conducted one of the first large-scale
  computational studies of ranks. For elliptic curves $y^2 = x^3 + ax + b$ with
  $|a|, |b| \leq B$, they found:
  \begin{center}
    \begin{tabular}{lr}
      \toprule
      Rank & Percentage \\
      \midrule
      0 & $\approx 35.2\%$ \\
      1 & $\approx 42.7\%$ \\
      2 & $\approx 18.9\%$ \\
      $\geq 3$ & $\approx 3.2\%$ \\
      \bottomrule
    \end{tabular}
  \end{center}
  These data are consistent with (but do not prove) Goldfeld's conjecture and the
  parity conjecture.
\end{remark}

% ==============================================================================
\section{2-Descent and Algorithms}\label{sec:descent}
% ==============================================================================

\begin{definition}[$2$-descent]
  A \emph{$2$-descent} on $E/\Q$ computes the $2$-Selmer group
  $\Sel^{(2)}(E/\Q)$ explicitly. Given $E\colon y^2 = (x-e_1)(x-e_2)(x-e_3)$
  (with $e_i \in \Q$ distinct), the $2$-Selmer elements correspond to
  quadratic twists $E_d$ with rational points (homogeneous spaces).
\end{definition}

\begin{algorithm}[$2$-descent algorithm, Birch-Swinnerton-Dyer]
  \begin{enumerate}[label=(\arabic*)]
    \item Compute $E[2](\Q)$ (the $2$-torsion).
    \item For each class $[c] \in (\Q^\times / (\Q^\times)^2)^2$ restricted to
      $S$-units (where $S$ = primes of bad reduction), check whether the corresponding
      homogeneous space has a rational point.
    \item Output $\Sel^{(2)}(E/\Q)$ as the classes with local points everywhere.
  \end{enumerate}
\end{algorithm}

\begin{remark}
  The $2$-Selmer rank gives an upper bound on $\rank E(\Q) + \dim_{\F_2} \Sha[2]$.
  To determine $\rank E(\Q)$ exactly, one needs to show $\Sha[2] = 0$ (not always provable).
\end{remark}

% ==============================================================================
\section{Open Problems}\label{sec:openproblems}
% ==============================================================================

The following are the main open problems concerning the rank of elliptic curves:

\begin{conjecture}[Birch-Swinnerton-Dyer, full version; Conjecture~\ref{conj:bsd}]
  Both parts of BSD remain open.
  The only known cases have analytic rank $\leq 1$.
\end{conjecture}

\begin{conjecture}[Goldfeld's conjecture; Conjecture~\ref{conj:goldfeld}]
  The average rank is $1/2$.
  Best known result: $\leq 0.885$ (Bhargava-Shankar).
\end{conjecture}

\begin{conjecture}[Boundedness conjecture; Conjecture~\ref{conj:bounded}]
  Is the rank uniformly bounded?
  Current record: $\geq 29$ (Elkies).
  All known constructions of high-rank curves are explicit; no theoretical barrier
  to arbitrarily high rank is known.
\end{conjecture}

\begin{conjecture}[Finiteness of $\Sha$; Conjecture~\ref{conj:sha}]
  Open in general; proven only for rank $\leq 1$ (conditionally).
\end{conjecture}

% ==============================================================================
\section{Conclusion}\label{sec:conclusion}
% ==============================================================================

The rank of elliptic curves is one of the most studied and most mysterious invariants
in number theory. The Mordell-Weil theorem establishes its finiteness, but the
determination of $r(E)$ for a given $E$ is in general a difficult open problem (no
unconditional algorithm is known). The BSD conjecture, if true, would give an
analytic handle on $r(E)$. Bhargava and Shankar's work has given the first statistical
control on the rank distribution, confirming that most elliptic curves have small rank.
The questions of whether the rank is bounded and what the average rank is remain
among the deepest unsolved problems in arithmetic geometry.

\begin{thebibliography}{99}

\bibitem{BS2015}
M.\ Bhargava, A.\ Shankar,
\textit{Binary quartic forms having bounded invariants, and the boundedness
of the average rank of elliptic curves},
Ann.\ of Math.\ (2) \textbf{181} (2015), no.\ 1, 191--242.

\bibitem{BS2015a}
M.\ Bhargava, A.\ Shankar,
\textit{Ternary cubic forms having bounded invariants, and the existence of a
positive proportion of elliptic curves having rank $0$},
Ann.\ of Math.\ (2) \textbf{181} (2015), no.\ 2, 587--621.

\bibitem{BS2013}
M.\ Bhargava, A.\ Shankar,
\textit{The average number of elements in the $4$-Selmer groups of elliptic curves is $7$},
J.\ Eur.\ Math.\ Soc.\ \textbf{18} (2016), no.\ 9, 2037--2089.

\bibitem{BSZ2014}
M.\ Bhargava, C.\ Skinner, W.\ Zhang,
\textit{A majority of elliptic curves over $\mathbb{Q}$ satisfy the Birch and
Swinnerton-Dyer conjecture},
preprint (2014), arXiv:1407.1826.

\bibitem{BM1990}
A.\ Brumer, O.\ McGuinness,
\textit{The behavior of the Mordell-Weil group of elliptic curves},
Bull.\ Amer.\ Math.\ Soc.\ \textbf{23} (1990), no.\ 2, 375--382.

\bibitem{BSD1965}
B.\,J.\ Birch, H.\,P.\,F.\ Swinnerton-Dyer,
\textit{Notes on elliptic curves. II},
J.\ Reine Angew.\ Math.\ \textbf{218} (1965), 79--108.

\bibitem{BCDT2001}
C.\ Breuil, B.\ Conrad, F.\ Diamond, R.\ Taylor,
\textit{On the modularity of elliptic curves over $\mathbf{Q}$: wild $3$-adic exercises},
J.\ Amer.\ Math.\ Soc.\ \textbf{14} (2001), no.\ 4, 843--939.

\bibitem{CW1977}
J.\ Coates, A.\ Wiles,
\textit{On the conjecture of Birch and Swinnerton-Dyer},
Invent.\ Math.\ \textbf{39} (1977), no.\ 3, 223--251.

\bibitem{Elkies2006}
N.\ Elkies,
\textit{$\mathbb{Z}^{28}$ in $E(\mathbb{Q})$, and $\mathbb{Z}^{29}$},
Postings to NMBRTHRY list (May 2006), available online.

\bibitem{Goldfeld1979}
D.\ Goldfeld,
\textit{Conjectures on elliptic curves over quadratic fields},
in \textit{Number Theory, Carbondale 1979}, Lect.\ Notes in Math.\ \textbf{751},
Springer (1979), 108--118.

\bibitem{GZ1986}
B.\,H.\ Gross, D.\,B.\ Zagier,
\textit{Heegner points and derivatives of $L$-series},
Invent.\ Math.\ \textbf{84} (1986), no.\ 2, 225--320.

\bibitem{Kolyvagin1988}
V.\,A.\ Kolyvagin,
\textit{Finiteness of $E(\mathbb{Q})$ and $\Sha(E,\mathbb{Q})$ for a subclass of
Weil curves},
Izv.\ Akad.\ Nauk SSSR Ser.\ Mat.\ \textbf{52} (1988), 522--540.

\bibitem{Mazur1977}
B.\ Mazur,
\textit{Modular curves and the Eisenstein ideal},
Publ.\ Math.\ IHÉS \textbf{47} (1977), 33--186.

\bibitem{Mestre1982}
J.-F.\ Mestre,
\textit{Construction d'une courbe elliptique de rang $\geq 12$},
C.\ R.\ Acad.\ Sci.\ Paris \textbf{295} (1982), 643--644.

\bibitem{Mordell1922}
L.\,J.\ Mordell,
\textit{On the rational solutions of the indeterminate equation of the third and fourth
degrees},
Proc.\ Cambridge Philos.\ Soc.\ \textbf{21} (1922), 179--192.

\bibitem{Nekovar2009}
J.\ Nekovář,
\textit{On the parity of ranks of Selmer groups IV},
Compos.\ Math.\ \textbf{145} (2009), no.\ 6, 1351--1359.

\bibitem{PPVW2019}
J.\ Park, B.\ Poonen, J.\ Voight, M.\,M.\ Wood,
\textit{A heuristic for boundedness of ranks of elliptic curves},
J.\ Eur.\ Math.\ Soc.\ \textbf{21} (2019), no.\ 9, 2859--2903.

\bibitem{PR2012}
B.\ Poonen, E.\ Rains,
\textit{Random maximal isotropic subspaces and Selmer groups},
J.\ Amer.\ Math.\ Soc.\ \textbf{25} (2012), no.\ 1, 245--269.

\bibitem{Rohrlich1993}
D.\ Rohrlich,
\textit{Variation of the root number in families of elliptic curves},
Compos.\ Math.\ \textbf{87} (1993), no.\ 2, 119--151.

\bibitem{Weil1928}
A.\ Weil,
\textit{L'arithmétique sur les courbes algébriques},
Acta Math.\ \textbf{52} (1928), 281--315.

\end{thebibliography}

\end{document}
